\section{Summarization: evidence, discovery and measured scoring cost}
\label{app:summeval-confirmation}
This experiment evaluates certificate size and discovery cost on summarization, using the 100-document, sixteen-system SummEval collection \citep{fabbri-etal-2021-summeval}. Before inspecting certification outcomes, we fixed a document-disjoint 33/33/34 reference/calibration/test split with seed 20261011. We also fixed four depths $q\in\{.05,.10,.20,.30\}$, $\alpha=.10$, both pooling rules and the primary and exact stopping criteria. Calibration documents are reserved and unused. The test batch contains 544 outputs. We report all prespecified acquisition configurations without filtering on query savings. This is a new task evaluation after method development, not a claim of SummEval's absence from evaluator training.

\subsection{Frozen cached-score confirmation}
Five metric families supply BARTScore CNN source-to-summary \citep{yuan2021bartscore}, BERTScore F \citep{zhang2020bertscore}, MoverScore \citep{zhao-etal-2019-moverscore}, Prism average \citep{thompson-post-2020-automatic} and ROUGE-L F \citep{lin-2004-rouge} scores. The cached matrix is obtained from \texttt{DAMO-NLP-SG/LLM\_summeval}, revision \texttt{bd7dde15817a} (full revision in the supplement). This experiment reuses released metric scores after identity and text-alignment checks; it does not rerun the original scorers. All 1,600 document/system identities align with the original SummEval annotations. Text alignment uses case, whitespace, punctuation and Penn Treebank bracket normalization; this does not establish identical scoring preprocessing. The raw source is retained locally with its provenance.

All ten choices of three base metrics and two candidates are evaluated. Higher scores are used as better by metric definition; human labels do not orient scores, select queries or choose stops. Each composition refits reference-only right ECDF maps, cutoff and ridge priorities. Ten compositions, four depths and two pooling rules give 80 instances and 160 certification targets. Seven acquisition policies give 560 complete paths and 1,120 stopping states. The integer budget is $\lfloor2kMT\alpha\rfloor$ with $M=3,T=2$. Independent replay checks 83,886 replies and every first stopping state.

Table~\ref{tab:summeval-strata} applies the same ex post capacity stratification used for the earlier NLP panels. The full-retention stratum has $V^*=|S_0|$; the exclusion stratum has $V^*+\varepsilon<|S_0|$; the remaining near-full stratum has $V^*<|S_0|\le V^*+\varepsilon$. These conditions are computed after acquisition and are not inputs to the policy. Under mean pooling, joint frontier saves two queries on average in the exclusion stratum at either stop. It increases the mean query count in the full-retention stratum. Under median pooling, joint frontier does not lower the average query count relative to safety-only in the exclusion stratum. Across the eight near-full mean primary targets, joint frontier uses 0.75 more queries on average: one target uses fewer, four tie and three use more than safety-only. Thus the mean-pooling saving in the exclusion stratum extends to this task conditionally; it is neither a universal advantage nor a per-instance guarantee. The compositions and depths share one test partition and are not independent replications.
\begin{table*}[t]
\centering\small
\begin{tabular}{lllrrrrr}\toprule
Pooling & Stop & Stratum & Targets & Less & Equal & More & Mean query difference\\\midrule
mean & exact & exclusion & 24 & 7 & 15 & 2 & -2.000 \\
mean & exact & full & 16 & 0 & 11 & 5 & 1.188 \\
mean & primary & exclusion & 16 & 7 & 7 & 2 & -2.000 \\
mean & primary & full & 16 & 0 & 10 & 6 & 1.875 \\
mean & primary & near full & 8 & 1 & 4 & 3 & 0.750 \\
median & exact & exclusion & 22 & 5 & 11 & 6 & 0.136 \\
median & exact & full & 18 & 1 & 12 & 5 & 1.222 \\
median & primary & exclusion & 18 & 4 & 10 & 4 & 0.167 \\
median & primary & full & 18 & 2 & 11 & 5 & 1.722 \\
median & primary & near full & 4 & 2 & 0 & 2 & 1.750 \\
\bottomrule\end{tabular}
\caption{Joint frontier versus safety-only acquisition on the 160 cached-score SummEval targets: ten evaluator compositions, four depths, two pooling rules and two stopping rules at $\alpha=.10$, with primary $\varepsilon=\lceil.02|S_0|\rceil$ and exact $\varepsilon=0$. The final column is joint-frontier minus safety-only entry queries. Less/equal/more count query differences, not certificate-minimum differences. All 160 targets are included; strata use full information only after acquisition.}
\label{tab:summeval-strata}\end{table*}


For the 48 non-full-retention mean targets, each joint and best separately optimal union program closes its own incumbent--lower-bound gap numerically. Reconstructed integer witnesses satisfy both count constraints. Joint evidence is smaller on 27 targets and equal on 21; numerical bound closure is not an exact rational minimality proof. Re-solving all 48 unions with HiGHS presolve disabled reproduced every value. The 68 full-retention targets across both pooling rules admit a direct integer computation described below. Odd-median targets also use the closed-form certificates, and enumeration of all four candidate-column subsets gives the minimum column count for $T=2$. Column enumeration does not establish entry-query minimality.

\paragraph{Exact computation in the full-retention subcase.}
If the true total selected-row cost is at most $b_H$, every partial lower-cost vector permits all selected rows. On $m$ active rows, the certificate therefore requires safety for at least $m-\min(m,\varepsilon)$ rows; exclusion is automatic. For row $i$, sort the reductions $p_i-f_{ir}$ in descending order and let $c_i(a)$ be $Tp_i$ minus the sum of the first $a$ reductions. Among all $a$ observations on this row, this is the smallest attainable upper cost. A dynamic program over processed rows, dropped-row count and total queries stores the minimum sum of retained upper costs. A retained row contributes $c_i(a)$ for each $a\in\{0,\ldots,T\}$; a dropped row contributes zero cost and zero queries. The smallest query count with at most $\varepsilon$ drops and cost at most $b_H$ is the exact minimum in the rectangular completion model.

To justify the recurrence, queries on dropped rows are unnecessary, and replacing a retained row's queried set by its largest reductions cannot worsen safety. Conversely, each dynamic-program transition specifies realizable observations and a retained subset. The recurrence therefore covers an optimal certificate and reconstructs a feasible one at its reported cost. Pairing the minimum safety witness with empty exclusion yields the same best separate union. This computation covers all 68 full-retention targets; exhaustive comparison on 60 small instances checks the implementation, not the general argument. It provides no characterization of arbitrary non-full-retention instances.

\subsection{Actual scorer execution}
\label{app:real-scoring}
To measure whether reduced score acquisition lowers execution time, a second experiment recomputes scores from the source texts on the same document split. It fixes BARTScore, BERTScore and MoverScore as base metrics and Prism and ROUGE-L as candidates. This is a fresh panel, not a reproduction of the cached matrix. BERTScore uses unrescaled RoBERTa-large layer-17 F with IDF disabled; MoverScore uses the original BERT-MNLI unigram power-mean implementation. Its two IDF dictionaries are fitted separately, on reference-split system summaries and on the human summaries of reference-split documents. Prism uses the bidirectional average of model m39v1 with a single reference. ROUGE-L uses \texttt{rouge-score} 0.1.2 with stemming and a single reference. BARTScore uses \texttt{facebook/bart-large-cnn} for source-to-summary likelihood. The scorer wrappers follow \texttt{neulab/BARTScore} revision \texttt{248f511cb34a} (full revision in the supplement); the supplement records exact model revisions and dependencies.

One A40 GPU executes FP32 scoring with TF32 disabled, PyTorch 2.4.1/CUDA 12.4 and four CPU threads. Each run loads models from local weights, scores all 528 reference outputs with five metrics and all 544 test outputs with three base metrics, fits maps and priorities, and then runs the specified candidate-scoring policy. Full scoring calls each metric once per scoring phase, processing inputs internally in batches of eight. Acquisition calls the candidate scorer once per requested entry without reading a precomputed candidate-test matrix. All methods include the common preparation, including the ridge fit. Eight methods, four depths, two pooling rules and three repetitions give 192 executions in a seeded random order. Each execution records the primary stop on its path to the exact stop, so the two stopping summaries reuse the same executions. Model downloads, dependency installation and post-path rescore audits are outside the measured pipeline interval.

All 192 runs have identical mapped reference and base arrays, despite small floating-point differences in raw scores. Every queried entry is rescored in a batch after the path finishes, and all mapped scores agree. Saved traces reconstruct every first stopping state. One early full-scoring run lacks the full candidate-array archive; its timing and counters remain in the complete grid, while the other 23 full runs retain these arrays. A score-entry count, a metric function invocation and an internal neural-network forward pass are distinct units. Tables~\ref{tab:real-scoring-primary} and~\ref{tab:real-scoring-exact} report the first two along with elapsed time; detailed forward counters and every run accompany the analysis.
\begin{table*}[t]\centering\small
\begin{tabular}{lrrrrr}\toprule
Method & Candidate entries & Metric calls & Total seconds & Time/full & Allocated USD\\\midrule
Full scoring & 1088.00 & 10.00 & 108.18 & 1.000 & 0.01797 \\
Reference & 120.25 & 128.25 & 89.97 & 0.832 & 0.01495 \\
Width & 123.25 & 131.25 & 90.43 & 0.836 & 0.01502 \\
Fixed & 125.25 & 133.25 & 89.51 & 0.827 & 0.01487 \\
Random & 124.12 & 132.12 & 89.95 & 0.832 & 0.01494 \\
Joint frontier & 123.50 & 131.50 & 90.24 & 0.834 & 0.01499 \\
Safety-only & 124.00 & 132.00 & 90.26 & 0.834 & 0.01499 \\
Exclusion-only & 122.88 & 130.88 & 90.37 & 0.835 & 0.01501 \\
\bottomrule\end{tabular}
\caption{Fresh SummEval scoring at $\alpha=.10$ and primary stopping, $V^+-V^-\le\lceil.02|S_0|\rceil$. Means cover mean and three-member median pooling, $q\in\{.05,.10,.20,.30\}$ and three repetitions (24 runs per method). Total time includes model loading and reference/base scoring. Time/full averages paired ratios. Metric calls include reference/base calls and may process multiple outputs; USD is duration times USD 0.598/hour, not a per-run invoice.}
\label{tab:real-scoring-primary}\end{table*}

\begin{table*}[t]\centering\small
\begin{tabular}{lrrrrr}\toprule
Method & Candidate entries & Metric calls & Total seconds & Time/full & Allocated USD\\\midrule
Full scoring & 1088.00 & 10.00 & 108.18 & 1.000 & 0.01797 \\
Reference & 125.88 & 133.88 & 90.19 & 0.834 & 0.01498 \\
Width & 128.88 & 136.88 & 90.48 & 0.836 & 0.01503 \\
Fixed & 131.00 & 139.00 & 89.73 & 0.829 & 0.01491 \\
Random & 129.25 & 137.25 & 90.20 & 0.834 & 0.01498 \\
Joint frontier & 129.88 & 137.88 & 90.49 & 0.837 & 0.01503 \\
Safety-only & 130.25 & 138.25 & 90.47 & 0.836 & 0.01503 \\
Exclusion-only & 128.50 & 136.50 & 90.45 & 0.836 & 0.01502 \\
\bottomrule\end{tabular}
\caption{Fresh SummEval scoring at $\alpha=.10$ and exact stopping, $V^+=V^-$. Means cover mean and three-member median pooling, $q\in\{.05,.10,.20,.30\}$ and three repetitions (24 runs per method). Total time includes model loading and reference/base scoring. Time/full averages paired ratios. Metric calls include reference/base calls and may process multiple outputs; USD is duration times USD 0.598/hour, not a per-run invoice.}
\label{tab:real-scoring-exact}\end{table*}


At primary stopping, reference priority uses 120.25 candidate-test entries per run on average, versus 1,088 for full scoring. Including the shared reference and base scoring, the mean score-entry totals are 4,392.25 and 5,360, so the total entry count falls far less than candidate-test entries. Mean total time falls from 108.18 to 89.97 seconds. At the same time, metric function calls rise from 10 to 128.25 because the on-demand interface uses single-entry calls. Each of the 24 paired reference runs is faster than its corresponding full run; this is a descriptive hardware-specific comparison, not a claim about all scorer deployments. Joint frontier and safety-only have similar total times on this fresh-score panel. Fewest candidate queries do not imply lowest mean time: fixed order averages 89.51 seconds against 89.97 for reference priority. This 0.46-second gap is smaller than every within-configuration range width for these two policies (Table~\ref{tab:runtime-variation}). The smaller offline joint certificates on some cached-score targets do not by themselves establish a runtime advantage between these policies.

Allocated compute cost multiplies each pipeline duration by the RunPod A40 GPU-plus-container quote of USD 0.598 per hour, recorded on 11 October 2026. Primary full and reference runs average USD 0.01797 and USD 0.01495, respectively. These are duration-based allocations, not observed per-run invoices or total account expenditure: installation, downloads, rescore audits, process startup and idle time also consume rented resources. The complete execution queue took 19,217.75 seconds. Table~\ref{tab:runtime-variation} reports within-configuration runtime variation; the supplement retains all 128 method/pooling/depth/stop summaries. Time ratios pair runs by pooling, depth and repetition. The comparison fixes batch-eight full scoring against sequential entry acquisition; online microbatching, persistent model services and alternative hardware may change the tradeoff.

\begin{table*}[t]\centering\small
\begin{tabular}{lrrrr}\toprule
Method & Primary SD range (s) & Exact SD range (s) & Primary range width (s) & Exact range width (s)\\\midrule
Fixed & 0.34--1.61 & 0.32--1.61 & 0.66--3.13 & 0.62--3.13 \\
Exclusion-only & 0.36--1.34 & 0.36--1.34 & 0.64--2.55 & 0.64--2.56 \\
Joint frontier & 0.19--3.05 & 0.19--3.06 & 0.37--5.33 & 0.37--5.33 \\
Safety-only & 0.26--1.97 & 0.26--1.97 & 0.48--3.44 & 0.49--3.44 \\
Full & 0.33--1.39 & 0.33--1.39 & 0.59--2.76 & 0.59--2.76 \\
Width & 0.32--1.29 & 0.32--1.29 & 0.60--2.30 & 0.60--2.30 \\
Random & 0.66--1.76 & 0.66--1.76 & 1.25--3.20 & 1.26--3.20 \\
Reference & 0.31--1.79 & 0.32--1.79 & 0.61--3.56 & 0.61--3.56 \\
\bottomrule\end{tabular}
\caption{Runtime variation within fixed scoring configurations. Each sample standard deviation (SD) and range width uses three repetitions at one pooling/depth setting. The table reports their minimum and maximum across the eight settings for each method; it does not pool configurations to estimate run noise. Complete per-setting means, ranges and SDs accompany the supplement.}
\label{tab:runtime-variation}\end{table*}

